\section{Discussion}
The studies establish complementary acceptance boundaries. Acquisition supplies qualified sources; frozen repair reuse converts development evidence into executable actions. The converter baseline shows that perfect audited topology can still omit physical values and supervision. R2G connects these layers through canonical identity, unit-aware extraction, stage visibility, and explicit validity. Downstream results close the evidence chain: geometry materially improves wirelength and resistance prediction, whereas hold slack is insensitive and feature-only models outperform GINE on several targets. Rather than claiming universal model dominance, the data support controlled tests of when topology and geometry matter.

\textbf{Reproducibility.} The fixed split, three-seed checkpoints, predictions, and per-design metrics are archived; a local rescore reproduces all 142,996 targets in a representative test run within cross-device tolerance.

\paragraph{Conclusion.} R2G connects qualified acquisition, frozen repair knowledge, and physically grounded graph construction through evidence contracts. It supplies 200 qualified sources, repairs eight of 13 held-out challenges, passes all core audits on 24 conversion designs, and produces stage-aware graphs with measurable physical signal: an auditable path from public RTL to aligned topology, geometry, and supervision.
